\Askhsh{1801}{4}
\begin{ylh}{1}
    \item 4.2. Τέμνουσα δύο ευθειών - Ευκλείδειο αίτημα
    \item 5.2. Παραλληλόγραμμα
    \item 5.6. Εφαρμογές στα τρίγωνα.
\end{ylh}
Δίνεται τρίγωνο $\mathrm{A}\mathrm{B}\Gamma$ με $\mathrm{A}\Gamma > \mathrm{A}\mathrm{B}$ και $\Delta, \mathrm{E}, \mathrm{Z}$ τα μέσα των πλευρών του $\mathrm{B}\Gamma, \mathrm{A}\Gamma, \mathrm{A}\mathrm{B}$ αντίστοιχα. Αν η διχοτόμος της γωνίας $\hat{\mathrm{B}}$ τέμνει την $\mathrm{Z}\mathrm{E}$ στο σημείο $\mathrm{M}$ και την προέκταση της $\Delta\mathrm{E}$ στο σημείο $\mathrm{N}$, να αποδείξετε ότι:
\begin{alist}
    \item Το τετράπλευρο $\mathrm{Z}\mathrm{E}\Delta\mathrm{B}$ είναι παραλληλόγραμμο. \monades{7}
    \item Τα τρίγωνα $\mathrm{B}\mathrm{Z}\mathrm{M}$ και $\mathrm{M}\mathrm{E}\mathrm{N}$ είναι ισοσκελή. \monades{10}
    \item $\mathrm{B}\mathrm{Z} + \mathrm{N}\mathrm{E} = \Delta\Gamma$. \monades{8}
\end{alist}
